% Transcription par lecture visuelle directe, sans OCR.
% Statut : transcrit ; relecture indépendante à effectuer.
% Chapitre complet, p. 315–331 (début) ; le chapitre second commence ensuite p. 331.
% Source : PDF original, vue 157, côté droit, page manuscrite 315.
\part*{Mémoires historiques pour servir à l'histoire de la Congrégation des hommes et des jeunes gens de Lyon de 1817 à 1840\\[1em]{\normalfont\large\textit{Ad majorem Dei Gloriam}}\\[0.5em]{\normalfont\small Pour la plus grande gloire de Dieu. (Note du transcripteur)}}
\addcontentsline{toc}{part}{Mémoires historiques, de 1817 à 1840}
\setcounter{chapter}{0}
\renewcommand{\theHchapter}{memoires.\arabic{chapter}}
\chapter{Fondation de la Congrégation des hommes}

\section*{Réflexions sur ces mémoires}
% Mention marginale de la page manuscrite 315 ; « Réflexions » répété p. 316 et 318.
Nous avons pu réussir, Messieurs, grâce à la bonté divine qui nous a visiblement secondés, à suivre dans l'histoire de la Congrégation que nous venons de terminer, et dans le précis historique dont elle est précédée, la marche de notre société et à vous la retracer jusque dans les plus petits détails. Après avoir jetté un coup d'œil soit sur les circonstances dont la Providence s'est servie pour préparer de loin ce pieux établissement, soit sur celles qu'elle a fait naître ensuite pour rapprocher et unir les premiers membres, nous avons considéré avec un tendre intérêt les foibles commencements de cette Congrégation naissante ; nous nous sommes sentis pénétrés d'une émotion bien vive en observant les effets de la tendresse maternelle avec laquelle notre bonne mère a pris soin de cette petite société lorsqu'elle étoit encore au berceau. Nous avons vu ensuite s'élever le vent de la persécution ; la Congrégation a courbé la tête ; elle s'est effacée pour laisser passer l'orage ; une main puissante l'a soutenue et elle n'a pas été renversée. Néanmoins le Seigneur, dans les conseils de sa sagesse qui sont aussi ceux de la miséricorde, a permis
% Source : PDF original, vue 158, côté gauche, page manuscrite 316.
pour la fortifier dans la foi, qu'elle trempa ses lèvres dans le calice d'amertume réservé au chef de l'Église. Deux de ses membres ont été jugés dignes de marcher à la suite du vicaire de J. C. dans la voie de la persécution. Enfin lorsqu'il a plu au souverain maître de manifester sa gloire, les ennemis de son saint nom ont été confondus et la Congrégation, disposée plus que jamais à se dévouer à la gloire de Dieu, a reparu triomphante à la suite de l'Église sa mère pour reprendre sa place dans les rangs de la milice chrétienne.

N'en doutons pas, c'est Marie, notre bonne mère qui a conservé dans ma mémoire tous ces souvenirs afin qu'il me fût possible de vous offrir un tableau fidèle des commencements de la Congrégation et des phases diverses de son existence pendant ces quinze premières années. Honneur en soit donc rendu à notre auguste patronne. Mais maintenant permettez-moi de vous communiquer ce qui se passe dans mon âme. Je dois vous le dire, je me sens effrayé de la tâche qui me reste à remplir ; voilà que la scène va singulièrement s'agrandir. Deux Congrégations ne faisant cependant qu'un seul tout vont marcher de pair ; leurs bonnes œuvres se multiplieront et se diversifieront à l'infini ; des sociétés nouvelles s'élèveront autour des nôtres. Directement ou indirectement la Congrégation mettra la main à tout, et se trouvera toujours pour quelque chose dans toutes les bonnes œuvres qui s'opéreront devant nos yeux. Comment au milieu de tant de détails parvenir à les saisir tous ? Comment surtout pouvoir bien me fixer pour garder l'ordre chronologique qui est cependant l'âme de l'histoire ?
% Source : PDF original, vue 158, côté droit, page manuscrite 317.
Plusieurs de ces difficultés me paroissent tout à fait insolubles. Dans le dénuement absolu où je suis de tout document, les obstacles s'élèvent devant mes yeux comme un rocher escarpé impossible à gravir ; aussi si j'étois seul je n'hésiterois pas à abandonner l'entreprise ; mais je compte encore sur l'assistance de notre bonne mère qui jusqu'ici a si visiblement conduit ma plume. Je me jette avec confiance dans ses bras, et fort de son secours, j'irai toujours en avant. Si je ne puis plus comme par le passé vous signaler à l'avenir les noms des préfets, si une multitude de circonstances intéressantes échappent à ma mémoire, j'ose espérer qu'elle m'en rappellera toujours quelques-unes, et que le bon Dieu permettra que ce soient celles qui seront les plus utiles à votre édification. D'ailleurs grâce aux procès verbaux qui tous ont été conservés depuis l'époque à laquelle nous sommes parvenus, et en s'aidant des notes qui ont été prises par les Préfets, il sera facile à celui d'entre nous qui voudra y consacrer quelques loisirs, en joignant les souvenirs à ceux que j'aurai conservés, de suppléer à ce qui manquera à cette relation et de rectifier les erreurs où le défaut de renseignements m'auroit entraîné.

Je trouve encore ici une occasion de remercier et de bénir la Providence, et d'admirer avec quelle bonté elle prend soin de nous jusque dans les moindres circonstances. Remarquez, Messieurs, qu'elle a laissé dans ma mémoire avec une entière clarté, jusqu'aux faits les plus minutieux de l'histoire de la Congrégation des premiers temps sur lesquels il n'y a rien d'écrit parce qu'elle vouloit que je pusse vous les transmettre avant de mourir ; mais c'est à l'époque précise où les documens
% Source : PDF original, vue 159, côté gauche, page manuscrite 318.
qui sont en votre pouvoir suffisent pour vous faire connoître ce que je ne pourrai pas vous retracer, que cette lumière qui m'a éclairé jusqu'ici semble s'affoiblir et presque m'abandonner.

J'espère toutefois, toujours avec le secours d'en haut, retrouver encore dans mes souvenirs une grande partie des particularités intéressantes de notre histoire qui, par leur nature, ne pouvoient laisser aucunes traces dans les procès verbaux. Continuons donc à suivre avec courage la carrière dans laquelle nous sommes entrés ; mais qu'il soit bien entendu que ce n'est plus une histoire mais seulement des mémoires qui vont se trouver sous ma plume.

\section*{Nécessité de fonder une Congrégation d'hommes}
% Mention marginale de la page manuscrite 318 ; date marginale : 1817.
La Congrégation agissoit en pleine liberté depuis près de trois ans. Elle s'étoit considérablement accrue, car vers le mois d'août 1817 soixante et dix membres environ prenoient sous la direction de M. Garnier qui étoit toujours Préfet, une part active à ses bonnes œuvres. En considérant le zèle de tous les Congréganistes et le dévouement avec lequel ils entreprenoient tout ce que leur société demandoit d'eux, on n'eût pas douté que la Congrégation n'eût alors tous les caractères de la vie et de la santé. Cependant au milieu de toute cette prospérité apparente les esprits attentifs apercevoient dans sa manière d'être un véritable germe de dissolution. Cette société existoit depuis quinze ans. Les Congréganistes des premières années étoient tous parvenus à l'âge mûr ; un grand nombre étoit marié ; un tiers de la totalité des membres au moins n'appartenoit plus à la classe des jeunes gens. Cette position devoit nécessairement s'aggraver chaque année par la force des choses ; on ne pouvoit donc pas se dissimuler que si on ne se hâtoit d'y apporter remède, bientôt il n'y auroit plus de Congrégation de jeunes gens, mais seulement une
% Source : PDF original, vue 159, côté droit, page manuscrite 319.
Congrégation bâtarde, une espèce de confrérie occupée de bonnes œuvres, où le mélange de tous les âges s'opposant à la conformité des goûts sans laquelle aucune intimité ne peut s'établir, il n'auroit jamais été possible de conserver un esprit propre. Outre ce grave inconvénient la présence des hommes au milieu des jeunes gens devoit nécessairement arrêter le développement des facultés de ceux-ci, et refroidir leur zèle. La timidité et la modestie, vertus si aimables dans la jeunesse et qui se rencontrent surtout dans les jeunes gens pieux, devoient naturellement les porter à s'effacer auprès de ces vétérans mûris par l'expérience et dont l'âge commandoit le respect. De là il fût évidemment résulté que tous les emplois et dignités auroient été constamment occupés par les anciens dont plusieurs auroient pu, même à 80 ans, faire encore partie de la Congrégation des jeunes gens, et la jeunesse se seroit vue réduite à composer la troupe des soldats sans jamais fournir un seul officier, dans une Congrégation formée pourtant tout exprès pour elle.

Un autre inconvénient se rencontroit encore dans la nature des œuvres ; toutes celles qui pouvoient être convenablement entreprises par des jeunes gens ne convenoient pas aussi bien à des pères de famille. La position que plusieurs d'entre eux occupoit dans le monde y mettoit obstacle, aussi quelques-uns avoient déjà commencé à se retirer, et il étoit à craindre que cet exemple ne fût suivi par d'autres ; il étoit donc urgent de trouver un moyen de remédier à un mal dont la Congrégation commençoit déjà à ressentir les funestes effets.

\section*{Congrégation des hommes}
% Mention marginale des p. 320, 324, 326 et 328 ; celle de p. 322 : « Fondation de la Congrégation des hommes ».
C'étoit déjà dans la prévision d'une situation qu'on regardoit encore comme éloignée, mais qui ne pouvoit pas ne pas se présenter
% Source : PDF original, vue 160, côté gauche, page manuscrite 320.
tôt ou tard qu'on avoit autrefois offert à la Congrégation de la sainte famille des arrangements qui auroient pu procurer à ces deux sociétés des avantages réciproques. Vous avez vu dans l'histoire de la Congrégation, que cette société avoit été emportée par la tempête. Depuis que la tranquillité étoit rétablie, quelques-uns de ses membres essayèrent plusieurs fois de provoquer une assemblée pour lui rendre la vie, mais tous leurs efforts furent inutiles ; ils ne purent jamais réussir à faire sortir le plus grand nombre de leurs frères de leur engourdissement, et bientôt tout espoir de restauration fut entièrement perdu.

La Congrégation des jeunes gens considéroit avec attention le résultat de ces démarches ; elle se proposoit si elles eussent été couronnées de quelque succès, d'entrer en négociation pour renouer s'il eût été possible le plan d'alliance dont jadis il avoit été question. Dès qu'il fut démontré qu'il n'y avoit plus rien à espérer de ce côté-là, le Préfet et ses assistants pensèrent que nous pourrions bien trouver en nous-même ce que nous avions vainement cherché à trouver en dehors. Il seroit difficile de ne pas apercevoir le doigt de Dieu en considérant l'inutilité des bons désirs de plusieurs anciens membres de la Congrégation des hommes, dans un moment où nous-même nous nous serions empressés d'autant plus d'aider cette Congrégation à se relever, que nous croyions alors y entrevoir pour nous de grands avantages. Ne vous paroît-il pas évident comme à moi, que si cette restauration a été impossible, c'est que dans les conseils de la sagesse éternelle, il avoit été résolu
% Source : PDF original, vue 160, côté droit, page manuscrite 321.
que la Congrégation des hommes sortiroit de la Congrégation des jeunes gens ; pour me servir d'une expression qui rend exactement ma pensée, elle devoit être la chair de sa chair et l'os de ses os, afin que ces deux sociétés n'eussent à elles deux qu'un seul esprit et qu'une seule âme.

Le préfet songea donc sérieusement à mettre ce grand projet à exécution. Il en conféra avec plusieurs anciens Congréganistes, surtout avec ceux que leur âge plus avancé rendoit immédiatement intéressés à cette mesure. À la suite de ces premières ouvertures et de l'avis du Père de la Congrégation, il choisit parmi ceux que leur âge devoit naturellement appeler à faire partie de la nouvelle Congrégation, cinq Congréganistes qu'il chargea de méditer sur les moyens d'arriver sans secousse à opérer la fondation désirée. Ces cinq Congréganistes après y avoir mûrement réfléchi, furent unanimes dans leur détermination. Ils arrêtèrent qu'on formeroit une Congrégation pour les hommes, mais qu'on la maintiendroit invariablement dans l'union la plus intime avec celle des jeunes gens, avec laquelle elle partageroit le titre de Congrégation de la Conception Immaculée, et que les deux Congrégations seroient toujours dirigées par le même père. Ces principes fondamentaux une fois établis, le Préfet désigna les cinq Congréganistes dont il avoit pris l'avis, pour être comme la première pierre de l'édifice qu'on se proposoit d'élever, et les réunit en commission pour préparer d'accord avec lui les règles de la Congrégation projettée.

Cette commission s'occupa aussitôt de son travail. Elle prit pour base du règlement qu'elle avoit à rédiger celui
% Source : PDF original, vue 161, côté gauche, page manuscrite 322.
de la Congrégation des jeunes gens. Tout ce qui concernoit le but, l'esprit et le gouvernement de la Congrégation fut conservé entièrement conforme. La nouvelle société fut divisée comme l'ancienne en trois sections ; seulement les œuvres et les dénominations furent en partie changées. La première section étoit celle de l'éducation ; elle devoit s'occuper de placer de jeunes enfants en apprentissage chez des maîtres religieux où ils pussent trouver de bons exemples et puiser de bons principes en même temps qu'ils apprendroient un état. La seconde étoit la section de charité ; sous un autre nom, elle devoit avoir les mêmes attributions que la section des aumônes. Enfin la troisième étoit la section de l'instruction. Elle devoit, tout en instruisant les ignorants, s'efforcer d'apporter un remède aux désordres dont gémissent la religion et la société, et de ramener les hérétiques au sein de la véritable Église. Tout en établissant pour la nouvelle Congrégation un comité de réception semblable à celui qui existoit dans l'ancienne, on en changea la dénomination ; il fut appelé comité du zèle, et non section du zèle, ce nom parut plus en rapport avec sa composition actuelle. Les assemblées générales ordinaires restèrent fixées à tous les mois, mais au Lundi en lieu du Dimanche. Les formes de réception n'éprouvèrent d'autres modifications que celles relatives à la Congrégation des jeunes gens dont nous parlerons plus bas. La Congrégation ne devoit recevoir que les personnes qui ne pouvoient pas entrer dans la Congrégation des jeunes gens, c'est-à-dire les hommes mariés et les garçons âgés de plus de trente-trois ans ; on limita à 60 ans l'âge où on pouvoit
% Source : PDF original, vue 161, côté droit, page manuscrite 323.
être reçu ; on ajouta aux prières journalières en usage dans la Congrégation des jeunes gens, un \textit{Pater} et un \textit{Ave} pour la conversion des pécheurs et un \textit{De profundis} pour les Congréganistes et leurs parents défunts\footnote{Débuts latins des prières : « Père », « Salut » et « Des profondeurs ». Traductions littérales. (Note du transcripteur).}. On fit quelques changements à la formule de la consécration à la Ste Vierge afin d'engager les Congréganistes, en se consacrant eux-mêmes à la mère de Dieu, à lui consacrer en même temps cette famille que la Providence a confiée à leur sollicitude.

Les fêtes que la Congrégation se proposa de célébrer en commun et par une communion générale furent à peu près les mêmes que celles déjà célébrées par l'ancienne ; il y eut cependant quelques différences. La Congrégation des hommes crut devoir placer au nombre de ses solennités la fête du Sacré Cœur. C'étoit un hommage solennel qu'elle fut heureuse de pouvoir rendre à ce cœur adorable. Elle savoit que dans ces derniers temps notre Seigneur sembloit avoir convié tous les fidèles à redoubler de confiance dans ce cœur si rempli d'amour pour les hommes, ce qui avoit été tant notre refuge dans les maux dont nous avions été environnés. Elle n'ignoroit pas que lorsque les malheurs engendrés par l'irréligion et que la Révolution a développés sont venus fondre sur notre France, notre saint et infortuné Roi Louis XVI avoit formé le projet de consacrer son royaume au cœur de Jésus et de solliciter du souverain Pontife l'établissement en son honneur d'une fête solennelle et chômée ; c'étoit donc pour la Congrégation s'associer autant qu'elle le pouvoit à ce pieux désir que d'inviter tous ses membres à venir chaque année au
% Source : PDF original, vue 162, côté gauche, page manuscrite 324.
jour de cette fête se plonger par une communion fréquente dans cet océan d'amour, afin d'obtenir la grâce d'aimer notre Dieu comme il mérite d'être aimé, ou du moins autant que nous sommes capables d'aimer.

La nouvelle Congrégation jugea aussi convenable de substituer dans son règlement la fête du Rosaire ou de N. D. des Victoires à celle de la Nativité de la Ste Vierge. C'est à cette dernière fête qu'étoit fixé, pour les jeunes gens, le commencement de l'année de la Congrégation et l'installation du Préfet. La commission dans la prévision que la Congrégation des hommes pourroit être fondée le jour de la fête du Rosaire fixa également à cette solennité l'époque annuelle de l'installation du Préfet des hommes. Il eût été plus difficile de choisir une époque plus convenable pour fonder la nouvelle Congrégation, car le premier sentiment que devoient éprouver les nouveaux membres de cette société naissante tous vieux Congréganistes de la Congrégation des jeunes gens, devoit être, en entrant dans cette nouvelle carrière, un sentiment bien vif de reconnoissance envers Marie par l'intervention de laquelle ils avoient déjà obtenu tant de faveurs. Cette société qu'ils alloient fonder ne devoit pas cesser d'être une partie inhérente de cette famille que notre bonne mère avoit adoptée depuis 15 ans. C'étoit encore sous ses auspices qu'elle se réunissoit. Oh ! comme leurs cœurs éprouvoient le besoin d'unir les saints transports de leur joie et de leur amour à ceux de tous les fidèles du monde catholique pour célébrer avec eux cette toute-puissance d'intercession dont le Seigneur du ciel et de la terre a revêtu celle
% Source : PDF original, vue 162, côté droit, page manuscrite 325.
qu'il a donnée pour mère à tous les chrétiens et pour refuge aux pauvres pécheurs, afin que les uns et les autres trouvassent constamment auprès de lui une avocate remplie de tendresse et de compassion, prête à demander et à obtenir pour eux toutes les grâces qui leur sont nécessaires !

La commission annexa au règlement qu'elle venoit de rédiger un projet de convention qui établissoit les rapports qui devoient exister entre les deux Congrégations. Cette convention imposoit aux jeunes gens l'obligation de passer dans la Congrégation des hommes lorsqu'ils seroient arrivés à un âge déterminé qui pouvoit cependant varier de quelques années, et à la Congrégation des hommes celle de les recevoir sans les soumettre à l'épreuve du scrutin. Elle réservoit aux Congrégations la faculté de célébrer, lorsqu'elles le jugeroient à propos, leurs fêtes en commun, et de se réunir en une seule assemblée dans quelques circonstances extraordinaires, et notamment chaque année pour entendre les rapports de toutes les sections. La convention imposoit encore aux Congrégations l'obligation d'être toujours dirigées par le même père, et étendoit à tous les membres des deux sociétés l'exercice des devoirs de charité que les règlements prescrivent aux Congréganistes entre eux.

Dès que ces deux pièces eurent été préparées, la commission jugea que ne se trouvant composée que de cinq membres, ce nombre n'étoit pas suffisant pour former, même au début, la nouvelle Congrégation. Elle crut donc convenable d'accord avec le préfet, d'en choisir quatre autres parmi les Congréganistes dont l'âge avoit marqué d'avance la place dans
% Source : PDF original, vue 163, côté gauche, page manuscrite 326.
la société nouvelle et elle se les adjoignit pour qu'ils vinssent avec ses membres concourir à la fondation de la Congrégation des hommes. Quelques jours avant le premier dimanche d'octobre ces neuf personnes se réunirent en une assemblée préliminaire ; le règlement et la convention furent de nouveau examinés avec grand soin, livrés à une discussion sérieuse et enfin définitivement adoptés à l'unanimité, sauf toutefois, quant à la convention, l'approbation de la Congrégation des jeunes gens, à laquelle on se proposoit de la soumettre seulement après la fondation. On prit ensuite toutes les mesures nécessaires à la conclusion de cette grande affaire ; on choisit, comme cela avoit été prévu d'avance, le jour de la fête du Rosaire pour celui de la fondation ; on convint qu'immédiatement après la messe on tiendroit une assemblée pour se constituer et élire un préfet, et que le soir même on communiqueroit à la Congrégation des jeunes gens qui devoit se réunir en assemblée générale, l'avis de la fondation de la Congrégation des hommes, et qu'on soumettroit à son approbation le projet de convention. Il fut d'ailleurs expressément convenu que les membres de la nouvelle Congrégation n'entendoient cesser de faire partie de la Congrégation des jeunes gens que lorsque la convention qui assuroit la continuation de l'union qui existoit entre eux et leurs anciens frères seroit acceptée par cette Congrégation. Toutes ces dispositions furent prises de concert avec le Préfet de la Congrégation des jeunes gens.

Tout étant ainsi déterminé les neuf fondateurs de la Congrégation des hommes, se rendirent le premier dimanche
% Source : PDF original, vue 163, côté droit, page manuscrite 327.
d'octobre, jour de la fête du saint Rosaire, dans la chapelle des Chartreux, où se trouva aussi M. Garnier, préfet de la Congrégation des jeunes gens. Je n'ose pas vous donner ici les noms des fondateurs, je craindrois de faire quelque erreur ; vous les trouverez inscrits dans les procès verbaux.

Voilà donc ces fondateurs de nouveau réunis au pied de cet autel de Marie, où déjà si souvent ils se sont rencontrés. C'est ici une circonstance qui distingue entièrement cette fondation de la première ; ces amis qui alloient contracter une nouvelle alliance se connoissoient depuis longtemps ; leur jeunesse s'étoit écoulée au service d'une mère commune ; c'étoit encore en se groupant autour d'elle qu'ils venoient aujourd'hui resserrer et rendre indissolubles jusqu'à la mort, ces nœuds que leurs cœurs pouvoient apprécier d'autant mieux que c'étoit depuis bien des années qu'ils en savouraient les douceurs ; ils avoient vieilli ensemble en marchant dans la même carrière en la présence de Dieu ; ils venoient aujourd'hui supplier le Seigneur de les maintenir tous ensemble dans cette voie qui seule peut conduire à la vie, jusqu'au moment suprême où l'éternité s'ouvrira devant eux. Jadis c'est en invoquant Marie qu'ils avoient consacré leur jeunesse au service de notre bon maître ; à cette heure c'est encore en invoquant Marie qu'ils viennent lui consacrer tout le temps qu'il leur reste à passer sur cette terre ; ils veulent aller toujours en avant, ne jamais regarder en arrière. C'est pourquoi ils renouvellent cette protestation déjà si souvent prononcée de renoncer à jamais à Satan, à ses pompes et à ses œuvres ; c'est à J. C. seul qu'ils veulent
% Source : PDF original, vue 164, côté gauche, page manuscrite 328.
appartenir et cela pour toujours, et malgré toutes les peines et les vicissitudes auxquelles ils pourront être condamnés. À cette résolution ferme et inébranlable, leurs cœurs, pénétrés de reconnoissance, ajoutent avec un vif empressement celle d'être constamment fidèles à la mère de Dieu ; aussi prononcent-ils avec bonheur cette consécration, que déjà au jour de leur entrée dans la Congrégation ils avoient faite de toute leur personne à notre auguste protectrice ; mais depuis le cercle de leurs affections s'est étendu, tout ce qui leur appartient doit appartenir à Marie ; c'est pourquoi, cette épouse, ces enfants chéris, ils les identifient avec eux ; avec eux, ils placent toute leur famille aux pieds de notre bonne mère et leurs yeux se mouillent de larmes en la lui consacrant. Marie est attentive à leurs prières ; elle les porte jusqu'au trône de l'Éternel, et dès lors plus d'un Congréganiste a pu s'apercevoir des bénédictions qu'elle a attirées sur tout ce qu'il a de plus cher. C'est en présence de J. C. que ces engagements solennels ont été prononcés et J. C. lui-même vient y mettre le sceau ; il vient habiter dans le cœur de tous, pour leur faire comprendre encore cette fois combien son joug est doux et son fardeau est léger.

\section*{Première assemblée. Élections}
% Mention marginale de la page manuscrite 328.
Après avoir cimenté leur union par la participation aux saints mystères les membres de la nouvelle Congrégation se formèrent en assemblée et nommèrent de suite, sans liste de candidats, ainsi que cela avoit été convenu pour cette première fois seulement le Préfet et les deux assistants. L'élection du secrétaire et du trésorier fut ajournée jusqu'à une époque où la Congrégation seroit
% Source : PDF original, vue 164, côté droit, page manuscrite 329.
devenue un peu plus nombreuse ; les assistants devoient jusque-là en remplir les fonctions.

Tous les membres de la Congrégation se rendirent avec exactitude à l'assemblée de la Congrégation des jeunes gens. Ils s'y présentèrent simplement comme membres de cette Congrégation, titre dont ils n'entendoient se départir qu'après l'acceptation de la convention. La fondation de la nouvelle Congrégation étoit encore complètement ignorée de la plupart des membres ; toutefois ils avoient pu la pressentir ; car le Préfet en publiant le jour de la fête de la Nativité après son installation, comme cela étoit l'usage alors, l'organisation de la Congrégation pour l'année courante, n'avoit donné aucun emploi aux cinq membres qui étoient chargés de préparer la fondation de la Congrégation des hommes, en annonçant qu'ils auroient une destination spéciale qu'il se réservoit de faire connoître plus tard.

Aussitôt après l'ouverture de l'assemblée et la lecture du procès verbal, le Préfet prit la parole ; il exposa la nécessité où se trouvoit la Congrégation de se fractionner, de former elle-même dans son sein une nouvelle Congrégation qui inséparablement unie avec l'ancienne, loin de rompre les liens qui existent entre nous ne feroit que les resserrer et les cimenter davantage. Il annonça en même temps que cette pensée dont il s'occupoit depuis longtemps, n'étoit plus seulement à l'état de projet, mais qu'elle avoit reçu le matin même un commencement d'exécution, que la nouvelle Congrégation venoit de se former au pied des autels, que déjà elle avoit élu son préfet ; que ce Préfet ainsi que tous les membres de la Congrégation nouvelle, étoient présents à l'assemblée et qu'ils avoient à lui
% Source : PDF original, vue 165, côté gauche, page manuscrite 330.
faire d'importantes communications. Après cette allocution le Préfet de la Congrégation des jeunes gens invita le Préfet de la Congrégation des hommes à venir prendre place au bureau. Celui-ci sortit alors des rangs des Congréganistes parmi lesquels il étoit confondu, et prit place à côté du Préfet et donna premièrement lecture du règlement adopté par la Congrégation des hommes, et ensuite du projet de convention qu'il soumit à l'approbation de l'assemblée, en en sollicitant l'acceptation.

Toutes ces révélations étonnèrent un peu la Congrégation ; il lui fut néanmoins facile de comprendre la sagesse de ces mesures et les avantages qu'elle y rencontreroit.

Un vote presque unanime vint donner la dernière sanction à tous ces arrangements ; je dis presque unanime, car comme il se trouva deux ou trois de nos amis, excellents Congréganistes sous tous les rapports et qui par l'âge qu'ils avoient atteint pouvoient prévoir qu'ils ne tarderoient pas beaucoup à passer dans la Congrégation des hommes, l'appréhension de cette séparation qu'ils regardoient comme un divorce, fascina pour lors si complètement les yeux de leur intelligence que la tendresse qu'ils avoient pour la Congrégation les engagea à s'opposer et à voter contre une mesure qui leur paroissoit désolante. Plus tard ils ont eux-mêmes reconnu que bien loin de vouloir séparer des amis, on ne s'étoit proposé au contraire autre chose que de rendre leurs cœurs de plus en plus inséparables, afin qu'ils fussent à jamais unis dans le cœur de Marie notre mère commune.

Les deux préfets apposèrent de suite leur signature à la
% Source : PDF original, vue 165, côté droit, page manuscrite 331 (début).
convention qui dès lors devint obligatoire pour l'une et l'autre Congrégation, et la Congrégation des hommes fut définitivement fondée.
